\chapter{Various}

\section{Intervals}
	\kactlimport{IntervalContainer.h}
	\kactlimport{IntervalCover.h}
	\kactlimport{ConstantIntervals.h}

\section{Misc. algorithms}
	\kactlimport{TernarySearch.h}
	\kactlimport{LIS.h}
	\kactlimport{FastKnapsack.h}

\section{Dynamic programming}
	\kactlimport{KnuthDP.h}
	\kactlimport{DivideAndConquerDP.h}

\section{Debugging tricks}
	\begin{itemize}
		\item \verb@signal(SIGSEGV, [](int) { _Exit(0); });@ converts segfaults into Wrong Answers.
			Similarly one can catch SIGABRT (assertion failures) and SIGFPE (zero divisions).
			\verb@_GLIBCXX_DEBUG@ failures generate SIGABRT (or SIGSEGV on gcc 5.4.0 apparently).
		\item \verb@feenableexcept(29);@ kills the program on NaNs (\texttt 1), 0-divs (\texttt 4), infinities (\texttt 8) and denormals (\texttt{16}).
	\end{itemize}

\section{Optimization tricks}
	\verb@__builtin_ia32_ldmxcsr(40896);@ disables denormals (which make floats 20x slower near their minimum value).
	\subsection{Bit hacks}
		\begin{itemize}
			\item \verb@x & -x@ is the least bit in \texttt{x}.
			\item \verb@for (int x = m; x; ) { --x &= m; ... }@ loops over all subset masks of \texttt{m} (except \texttt{m} itself).
			\item \verb@c = x&-x, r = x+c; (((r^x) >> 2)/c) | r@ is the next number after \texttt{x} with the same number of bits set.
			\item \verb@rep(b,0,K) rep(i,0,(1 << K))@ \\ \verb@  if (i & 1 << b) D[i] += D[i^(1 << b)];@ computes all sums of subsets.
		\end{itemize}
	\subsection{Pragmas}
		\begin{itemize}
			\item \lstinline{#pragma GCC optimize ("Ofast")} will make GCC auto-vectorize loops and optimizes floating points better.
			\item \lstinline{#pragma GCC target ("avx2")} can double performance of vectorized code, but causes crashes on old machines.
			\item \lstinline{#pragma GCC optimize ("trapv")} kills the program on integer overflows (but is really slow).
		\end{itemize}
	\kactlimport{FastMod.h}
	\kactlimport{FastInput.h}
	\kactlimport{BumpAllocator.h}
	\kactlimport{SmallPtr.h}
	\kactlimport{BumpAllocatorSTL.h}
	% \kactlimport{Unrolling.h}
	\kactlimport{SIMD.h}
\section{Information about numbers}
			How to read this table: [Upper bound] | [Max. number of divisors] | [Max. number of distinct prime divisors] | [Smallest primes greater than bound] | [Max. difference between primes] | [Sum of Divisor Count]
			
			\begin{tabular}{|c|c|c|c|c|c|} \hline
				$10^2$ & 12 & 3 & 1, 3, 7, 9, 13 & 8                & 473 \\ \hline
				$10^3$ & 32 & 4 & 9, 13, 19, 21, 31 & 20            & 7053 \\ \hline
				$10^4$ & 64 & 5 & 7, 9, 37, 39, 61 & 36             & 93643 \\ \hline
				$10^5$ & 128 & 6 & 3, 19, 43, 49, 57 & 72           & 1166714 \\ \hline
				$10^6$ & 240 & 7 & 3, 33, 37, 39, 81 & 114          & 13969985 \\ \hline
				$10^7$ & 448 & 8 & 19, 79, 103, 121, 139 & 154      & 162725300 \\ \hline
				$10^8$ & 768  & 8 & 7, 37, 39, 49, 73 & 220         & 1857511487 \\ \hline
				$10^9$ & 1344 & 9 & 7, 9, 21, 33, 87 & 282          & 20877697534 \\ \hline
				$10^{10}$ & 2304 & 10 & 19, 33, 61, 69, 97 & 354    & - \\ \hline
				$10^{12}$ & 6720 & 11 & 39, 61, 63, 91, 121 & 540   & - \\ \hline
				$10^{14}$ & 17280 & 13 & 31, 67, 97, 99, 133 & 804  & - \\ \hline
				$10^{16}$ & 41472 & 14 & 61, 69, 79, 99, 453 & 1132 & - \\ \hline
				$10^{18}$ & 103680 & 15 & 3, 9, 31, 79, 177 & 1442  & - \\ \hline
			\end{tabular}

\section{C++ Math functions}
	\begin{tabular}{|c|c|} \hline
		acos & [-1, 1] $\rightarrow$ [0, $\pi$] \\ \hline
		asin & [-1, 1] $\rightarrow$ [$-\frac{\pi}{2}$, $\frac{\pi}{2}$] \\ \hline
		atan(x) & $x \rightarrow$ [$-\frac{\pi}{2}$, $\frac{\pi}{2}$] \\ \hline
		atan2(y, x) & $\frac{y}{x} \rightarrow$ [$-\pi$, $\pi$] \\ \hline
	\end{tabular}
\section{Bitwise built-in functions}
	Remember: If you're using 64-bit integers, remember to use the l/ll version of the functions
	\begin{itemize}
  		\item \_\_builtin\_ffs: Returns one plus the index of the least significant 1-bit of x, or if x is zero, returns zero.
  		\item \_\_builtin\_clz: Returns the number of leading 0-bits in x, starting at the most significant bit position. If x is 0, the result is undefined.
  		\item \_\_builtin\_ctz: Returns the number of trailing 0-bits in x, starting at the least significant bit position. If x is 0, the result is undefined.
		\item \_\_builtin\_clrsb: Returns the number of leading redundant sign bits in x, i.e. the number of bits following the most significant bit that are identical to it.
		\item \_\_builtin\_popcount: Returns the number of 1-bits in x.
		\item \_\_builtin\_parity: Returns the parity of x, i.e. the number of 1-bits in x modulo 2.
	\end{itemize}
\section{Debug Template Niklas}
	\kactlimport{NiklasDebug.h}

\section{Lazy Segtree}
	\kactlimport{NiklasSegtree.h}

\section{Grid Hasher}
	\kactlimport{GridHasher.h}